% Teil II: durchgerechneter Weg, Teile a--c.

\subsection*{Gegeben und gesucht}
$a)$ Unterschied zwischen $x^2+2=0$ und $-D\,x(t)=m\,\ddot x(t)$; $b)$
$v(t),\ a(t)$ für $x(t)=x_0\cos(\omega t+\varphi_0)$; $c)$ Zeigen:
$\ddot x=-\omega^2x$, also Lösung der DGL $\ddot x+\omega^2x=0$.

\subsection*{Ansatz}
Beide Faelle unterscheiden, dann die trigonometrischen Ableitungen bilden
und die DGL prüfen.

\subsection*{Rechnung}
\textbf{a) Unterschied.} $x^2+2=0$ ist eine algebraische Gleichung; gesucht
sind die \emph{Zahlenwerte} $x$ (hier komplex $\pm i\sqrt2$), die die
Gleichung erfuellen. $-D\,x(t)=m\,\ddot x(t)$ ist eine
Differentialgleichung; gesucht ist die \emph{Funktion} $x(t)$, die die
Gleichung zusammen mit ihrer Ableitung identisch erfuellt.

\textbf{b) Ableitungen.}
\begin{align}
  v(t)&=\dot x(t)=-x_0\omega\sin(\omega t+\varphi_0),\\
  a(t)&=\ddot x(t)=-x_0\omega^2\cos(\omega t+\varphi_0)\;.
\end{align}

\textbf{c) Nachweis.} Wegen $x(t)=x_0\cos(\omega t+\varphi_0)$ ist
\begin{equation}
  a(t)=-x_0\omega^2\cos(\omega t+\varphi_0)=-\omega^2\,x(t)\;.
\end{equation}
Damit folgt
\begin{equation}
  \ddot x+\omega^2x=-\omega^2x+\omega^2x=0\;\checkmark.
\end{equation}

\begin{ergebnis}
$a)$ algebraische Gleichung: gesucht die Zahl; DGL: gesucht die Funktion.
$b)$ $v=-x_0\omega\sin(\omega t+\varphi_0)$, $a=-x_0\omega^2\cos(\omega t+\varphi_0)$.
$c)$ $a=-\omega^2x$, d. h. $x(t)$ loest $\ddot x+\omega^2x=0$.
\end{ergebnis}

\begin{kontrolle}
\textbf{Einheiten:} $\omega$ in \unit{\per\second}, $\omega t$ und
$\varphi_0$ dimensionslos; $x_0\omega^2$ hat die Dimension von $a$
(\unit{\metre\per\second\squared}), $x_0\omega$ die von $v$
(\unit{\metre\per\second}). \checkmark

\textbf{Querprobe:} Für $\varphi_0=0$, $t=0$ ist $x=x_0$ (maximal),
$v=0$ (Umkehrpunkt), $a=-x_0\omega^2$ nach innen -- klassischer
Oszillator-Umkehrpunkt. \checkmark
\end{kontrolle}

\begin{merksatz}
Der harmonische Oszillator beschreibt jede Bewegung mit rücktreibender
Kraft proportional zur Auslenkung: $x(t)=x_0\cos(\omega t+\varphi_0)$ ist
seine allgemeine Lösung und $\ddot x+\omega^2x=0$ seine Bewegungsgleichung.
\end{merksatz}
